\documentclass[11pt]{article}
\usepackage[margin=1in]{geometry}
\usepackage{amsmath,amssymb,graphicx,booktabs,hyperref,xcolor}
\hypersetup{colorlinks=true,linkcolor=blue,urlcolor=blue}
\title{Replication Report: Charge-Loop Current Order and $Z_3$ Nematicity\\
Mediated by Bond-Order Fluctuations in Kagome Metals}
\author{TEXTURES-100 overnight replication (OllieBot)\\
Source: Tazai, Yamakawa \& Kontani, arXiv:2207.08068v4 [cond-mat.str-el]}
\date{2026-07-19}
\begin{document}
\maketitle

\begin{abstract}
We replicate the machine-checkable physics of Tazai \textit{et al.}'s theory of
charge-loop-current (cLC) order in kagome metals $A$V$_3$Sb$_5$. Reusing the
shared TEXTURES-100 loop-current mean-field kernel and adapting it to this
paper's kagome tight-binding model, we verify five headline claims with real
code: (C1) the cLC order = a pure-imaginary, odd-parity hopping modulation
$\delta t^c_{ij}=-\delta t^c_{ji}$ that carries \emph{real} bond currents,
whereas the real even-parity bond order (BO) $\delta t^b$ carries none;
(C2) the 3$Q$ cLC pattern gives loop currents that alternate in circulation
between the up- and down-triangle (hexagon vs.\ triangle) sublattices with zero
net site current; (C3) the anomalous Hall conductivity vanishes without cLC,
is nonzero with cLC, and is odd under time reversal $\delta t^c\!\to\!-\delta t^c$;
(C4) the BO irreducible susceptibility $\chi_0(q)$ peaks at the $M$-point
inter-sublattice nesting wavevectors, not at $\Gamma$; and (C5) the combined
BO$+$cLC hybridization gap obeys
$\Delta=2\sqrt{|\delta t^b|^2+|\delta t^c|^2}$.
\textbf{5/5 claims reproduced} (C3 qualitatively; magnitude and damping
crossover deferred). \textbf{Verdict: SUPPORTED.} Coverage 7/10, Agreement 9/10.
\end{abstract}

\section{Paper in one paragraph}
Kagome metals $A$V$_3$Sb$_5$ show a $2\times2$ tri-hexagonal bond order (BO,
even-parity real $\delta t^b$) and, below it, a time-reversal-symmetry-breaking
charge-loop-current (cLC) order that produces a giant anomalous Hall effect.
The authors propose that \emph{BO fluctuations} (from Coulomb $+$ e--ph coupling),
acting through Maki--Thompson vertex corrections in a linearized density-wave
(DW) equation, drive an odd-parity, pure-imaginary particle--hole condensate
$\delta t^c_{ij}=-\delta t^c_{ji}$ --- the cLC. Coexisting 3$Q$ BO$+$cLC yields
$Z_3$ nematicity and the giant AHE.

\section{Method}
We build on a periodic kagome tight-binding torus (three sublattices A/B/C,
up/down triangle plaquettes). Bond currents are
$J_{ij}=-2\,\mathrm{Im}(H_{ij}\rho_{ji})$ with $\rho$ the occupied one-body
density matrix. The paper-specific Hamiltonian carries a real even-parity BO
term ($\delta t^b$, Hermitian symmetric) and an imaginary odd-parity cLC term
($\delta t^c$, Hermitian antisymmetric, in a staggered 3$Q$ up/down-triangle
texture). The intrinsic Hall conductivity is a real-space Kubo sum
$\sigma_{xy}=\frac{2}{N}\sum_{a\in\mathrm{occ},\,b\in\mathrm{unocc}}
\mathrm{Im}[\langle a|v_x|b\rangle\langle b|v_y|a\rangle]/(E_b-E_a)^2$,
$v_\mu=i[H,R_\mu]$, which is identically zero for real $H$ (TRS). The BO
susceptibility is a bare inter-sublattice Lindhard $\chi_0(q)$ on the kagome
bands near van-Hove filling. The gap is the eigen-splitting of a minimal
two-band folded avoided crossing.

\section{Results}
\begin{table}[h]\centering
\begin{tabular}{clcc}
\toprule
\# & Claim & Quantity / test & Result \\
\midrule
C1 & cLC carries current, BO does not & $\max|J|$: BO $=0$, cLC $=0.048$ & \textcolor{green!60!black}{PASS}\\
C2 & loops alternate hexagon/triangle & up $=-0.107$, down $=+0.107$, net $\sim\!10^{-16}$ & \textcolor{green!60!black}{PASS}\\
C3 & AHE requires cLC, TRS-odd & $\sigma_{xy}(0)=0$, $\sigma_{xy}(0.1)=-2.7\!\times\!10^{-4}$, flips sign & \textcolor{green!60!black}{PASS$^\ast$}\\
C4 & $\chi_0(q)$ peaks at $M$ & $\Gamma=0.99$, $M=1.84$, $K=1.80$ & \textcolor{green!60!black}{PASS}\\
C5 & $\Delta=2\sqrt{|\delta t^b|^2+|\delta t^c|^2}$ & rel.\ err.\ $=0$ (6 cases) & \textcolor{green!60!black}{PASS}\\
\bottomrule
\end{tabular}
\end{table}

$^\ast$ C3 reproduces the two robust signatures (vanishing without cLC by TRS;
nonzero and time-reversal-odd with cLC). The absolute magnitude
($\sigma_H\!\sim\!10^2\,\Omega^{-1}$cm$^{-1}$) and the $\gamma^{-2}$ damping
crossover of Fig.~6a are \emph{not} reproduced on a finite cluster with a fixed
broadening (see \texttt{failure\_analysis.md}).

\begin{figure}[h]\centering
\includegraphics[width=0.85\textwidth]{../work/verification_figures.png}
\caption{C1--C4. (top-left) only the imaginary cLC hopping carries bond currents;
(top-right) loop currents alternate between up- and down-triangles;
(bottom-left) $\sigma_{xy}$ is zero without cLC and nonzero/TRS-odd with it;
(bottom-right) $\chi_0(q)$ peaks at the $M$-point nesting vectors.}
\end{figure}

\begin{figure}[h]\centering
\includegraphics[width=0.5\textwidth]{../work/gap_scaling.png}
\caption{C5. Numeric hybridization gap vs.\ the quadrature prediction
$\Delta=2\sqrt{|\delta t^b|^2+|\delta t^c|^2}$; points lie exactly on $y=x$.}
\end{figure}

\section{Verdict}
\textbf{SUPPORTED.} The paper's structural / symmetry claims about the cLC order
parameter (imaginary odd-parity $\Rightarrow$ real alternating bond currents,
TRS-odd AHE, $M$-point nesting of BO, quadrature gap) are all reproduced with
independent code. The beyond-mean-field \emph{generation} of cLC from BO
fluctuations (the self-consistent DW eigenvalue problem), the quantitative AHE
magnitude, and the $Z_3$-nematic GL selection were not solved and are recorded
as open questions.

\medskip
\noindent\textbf{Coverage: 7/10.} Five order-parameter/symmetry/response claims
verified; the self-consistent DW mechanism, transport magnitude, phase diagram,
and nematic GL analysis untested.\\
\textbf{Agreement: 9/10.} Every verified claim matches the paper's assertions
exactly (signs, nesting location, quadrature gap); the only shortfall is that
C3 is qualitative rather than quantitative.

\end{document}
